\documentclass[11pt,a4paper]{article}

\usepackage{fontspec}
\usepackage[french]{babel}
\usepackage[margin=2cm]{geometry}
\usepackage{xcolor}
\usepackage{listings}
\usepackage{booktabs}
\usepackage{tabularx}
\usepackage{enumitem}
\usepackage{hyperref}

\definecolor{bleu}{RGB}{25,60,110}
\definecolor{fond}{RGB}{245,246,248}
\definecolor{vert}{RGB}{30,110,60}
\definecolor{gris}{RGB}{110,110,110}

\hypersetup{colorlinks=true, linkcolor=bleu, urlcolor=bleu}
\setlist{nosep, leftmargin=*}
\setlength{\parindent}{0pt}
\setlength{\parskip}{4pt}

\lstset{
  basicstyle=\ttfamily\small,
  backgroundcolor=\color{fond},
  frame=single, rulecolor=\color{gris!40},
  commentstyle=\color{vert},
  morecomment=[l]{\#},
  columns=fullflexible, keepspaces=true,
  breaklines=true, xleftmargin=4pt, xrightmargin=4pt
}

\newcommand{\retenir}[1]{\par\smallskip\noindent\colorbox{fond}{\parbox{\dimexpr\linewidth-2\fboxsep}{\textbf{\textcolor{bleu}{À retenir :}} #1}}\par\smallskip}

\title{\textbf{\textcolor{bleu}{Kubernetes : mon résumé de révision}}\\[4pt]
\large Des bases au rolling update et au rollback}
\author{Hicham Habchane}
\date{19 octobre 2026 --- Programme DevOps, Jour 1}

\begin{document}
\maketitle

\section{L'idée de Kubernetes}

\textbf{Docker} crée et lance des conteneurs. \textbf{Kubernetes} les organise : il les lance, les multiplie, les surveille et les relance.

\textbf{L'état souhaité :} je ne donne pas d'ordre (« lance un conteneur »), je \emph{décris} ce que je veux (« toujours 3 nginx »). Kubernetes corrige la réalité en permanence.

\begin{center}
\begin{tabular}{@{}ll@{}}
\toprule
\textbf{Kubernetes} & \textbf{Restaurant} \\
\midrule
Conteneur & le plat \\
Pod & le plateau \\
ReplicaSet & le chef d'équipe qui compte les plateaux \\
Deployment & le responsable de salle \\
template & la fiche recette \\
labels & l'autocollant sur le plateau \\
Service & l'accueil du restaurant \\
Kubernetes / kubectl & le gérant / la télécommande \\
\bottomrule
\end{tabular}
\end{center}

\section{Le cluster (minikube)}

Un seul nœud \texttt{minikube} : il dirige (control-plane) et fait tourner les Pods. Les conteneurs tournent avec \textbf{containerd}.

\begin{lstlisting}
minikube start        # allumer le cluster (après chaque redémarrage du PC)
minikube status       # allumé ou éteint ?
kubectl get nodes     # attendre STATUS = Ready
\end{lstlisting}
\retenir{erreur \texttt{no route to host} ou \texttt{connection refused} $\rightarrow$ le cluster est éteint.}

\section{Le Pod}

La plus petite unité de Kubernetes : une enveloppe autour d'un conteneur, avec un nom et une IP interne.

\begin{lstlisting}
kubectl apply -f pod-nginx.yaml     # créer le Pod
kubectl get pods                    # ContainerCreating puis Running
kubectl describe pod <nom>          # fiche détaillée, Events en bas
kubectl logs <nom>                  # ce que le conteneur affiche
\end{lstlisting}

Un fichier YAML a toujours 4 parties : \texttt{apiVersion}, \texttt{kind}, \texttt{metadata}, \texttt{spec}.

\retenir{un Pod est \textbf{jetable}. Un Pod seul supprimé ne revient \textbf{jamais}.\\
\texttt{RESTARTS} = redémarrages du \emph{conteneur}, pas du Pod.}

\section{Le ReplicaSet}

Il garde \textbf{N Pods identiques} en marche. Il ne regarde \textbf{que les étiquettes} (\texttt{selector}).

\begin{itemize}
  \item Pod \textbf{en trop} avec la bonne étiquette $\rightarrow$ il en \textbf{supprime} un (le Pod \texttt{intrus}).
  \item Pod orphelin déjà présent avec la bonne étiquette $\rightarrow$ il l'\textbf{adopte}.
  \item \textbf{Limite :} si je change l'image, les Pods existants \textbf{ne changent pas}.
\end{itemize}

\begin{lstlisting}
kubectl get pods -l app=web          # filtrer par étiquette
kubectl describe pod intrus          # Controlled By: ReplicaSet/web-rs = adoption
\end{lstlisting}
\retenir{1 sélecteur = 1 seul contrôleur. On n'écrit jamais de ReplicaSet à la main : on écrit un Deployment.}

\section{Le Deployment}

Il crée et pilote les ReplicaSets. Nom d'un Pod : \texttt{monsite-86f69f946d-4vcxz}
= Deployment -- \emph{hash du template} (le ReplicaSet) -- identifiant du Pod.

\begin{center}
\begin{tabularx}{\linewidth}{@{}lX@{}}
\toprule
\textbf{Je change\dots} & \textbf{Résultat} \\
\midrule
\texttt{replicas} (3 $\rightarrow$ 5) & même ReplicaSet, seul DESIRED change, anciens Pods gardés \\
le \texttt{template} (l'image) & \textbf{nouveau} ReplicaSet = nouvelle \textbf{révision} \\
\bottomrule
\end{tabularx}
\end{center}

\retenir{1 révision = 1 ReplicaSet. Les anciens ReplicaSets restent à 0 pour permettre le rollback.}

\section{Rolling update (mise à jour sans coupure)}

\begin{lstlisting}
sed -i 's/image: nginx$/image: nginx:1.27/' deployment-nginx.yaml
kubectl apply -f deployment-nginx.yaml
kubectl rollout status deployment monsite     # suivre en direct
\end{lstlisting}

Le nouveau ReplicaSet monte (0$\rightarrow$5) pendant que l'ancien descend (5$\rightarrow$0), petit à petit.

\begin{itemize}
  \item \textbf{maxSurge 25\,\%} (arrondi au-dessus) : jusqu'à 2 Pods \emph{en plus} (7 au total).
  \item \textbf{maxUnavailable 25\,\%} (arrondi au-dessous) : toujours au moins 4 Pods disponibles.
\end{itemize}
\retenir{pendant la mise à jour, 2 versions coexistent : elles doivent rester compatibles.}

\section{Rollback (retour arrière)}

\begin{lstlisting}
kubectl rollout history deployment monsite    # révisions
kubectl rollout undo deployment monsite       # revenir à la version d'avant
\end{lstlisting}

Le Deployment \textbf{réutilise l'ancien ReplicaSet} avec des Pods neufs. C'est un nouveau déploiement : la révision 1 devient la révision 3.

\section{Déclaratif, impératif et dérive}

\begin{itemize}
  \item \textbf{Déclaratif} (bonne pratique) : je modifie le fichier puis \texttt{apply}. Le fichier dans Git fait foi.
  \item \textbf{Impératif} : \texttt{kubectl scale}, \texttt{rollout undo}. Le fichier n'est pas modifié.
  \item Cluster $\neq$ fichier = \textbf{dérive} (drift) : le prochain \texttt{apply} remet l'ancien état.
\end{itemize}
\retenir{après un rollback, je remets le fichier d'accord avec le cluster.
\texttt{apply} répond \texttt{created}, \texttt{configured} ou \texttt{unchanged}.}

\section{Services (à venir)}

Les Pods changent d'IP. Le Service donne \textbf{un nom et une adresse fixe} devant un groupe de Pods (par \texttt{selector}) et répartit les requêtes.
\textbf{ClusterIP} : interne au cluster. \textbf{NodePort} : accès externe (ports 30000--32767). \textbf{port-forward} : tunnel temporaire vers mon PC.

\section{Questions d'entretien}

\begin{tabularx}{\linewidth}{@{}>{\bfseries}p{5.2cm}X@{}}
\toprule
Question & Réponse courte \\
\midrule
Pod, ReplicaSet, Deployment ? & Pod = exécute ; ReplicaSet = maintient N Pods par étiquettes ; Deployment = pilote les ReplicaSets, rolling update et rollback. \\
Comment K8s garde 3 replicas ? & Le ReplicaSet compare sans arrêt le nombre de Pods étiquetés au nombre voulu, puis crée ou supprime. \\
Pourquoi garder les vieux ReplicaSets ? & Pour le rollback. \\
maxSurge / maxUnavailable ? & Pods en plus autorisés / Pods indisponibles au maximum $\rightarrow$ pas de coupure. \\
Déclaratif ou impératif ? & Déclaratif : le fichier dans Git fait foi, sinon dérive. \\
\bottomrule
\end{tabularx}

\end{document}
